\par\indent\SourceLeaf{208}%
Setzt man für \(f\) z. B. eine Exponentialgrösse, nemlich
\[
f\upsilon=he^{\upsilon}
\]
\noindent wo \(h\) constant ist und \(e\) die Basis der hyperbolischen Logarithmen bedeutet, so hat man die einfachste Darstellung
\[
X=ht\cos\BareRoot{\frac{kk}{kk+1}}\mathord{.}u,\quad
Y=ht\sin\BareRoot{\frac{kk}{kk+1}}\mathord{.}u
\]

Die Anwendung der Formeln des 7. Art. gibt hier
\[
n=(kk+1)tt,\quad N=1
\]
\noindent und, da \(\varphi\upsilon=\varphi'\upsilon=he^{\upsilon}\),
\[
\varphi(\log t+i\BareRoot{\frac{kk}{kk+1}}\mathord{.}u)\mathord{.}\varphi'(\log t-i\BareRoot{\frac{kk'}{kk+1}}\mathord{.}u)=hhtt
\]
\noindent folglich
\[
m=\frac{h}{\BareRoot{(kk+1)}}
\]
\noindent also constant. Macht man also noch
\[
h=\BareRoot{(kk+1)}
\]
\noindent so wird die Darstellung eine vollkommne Abwicklung.

\Article{10}
Es sei drittens die Kugelfläche, deren Halbmesser \(=a\), in der Ebne darzustellen. Wir setzen hier
\[
\begin{aligned}
x&=a\cos t\mathord{.}\sin u\\
y&=a\sin t\mathord{.}\sin u\\
z&=a\cos u
\end{aligned}
\]
\noindent wodurch wir erhalten
\[
\omega=aa\sin u^2\,\mathrm{d}t^2+aa\,\mathrm{d}u^2
\]
\noindent Die Differentialformel \(\omega=0\) gibt folglich
\[
\mathrm{d}t\mp i\mathord{.}\frac{\mathrm{d}u}{\sin u}=0
\]
\noindent und deren Integration
\[
t\pm i\log\mathrm{cotang}\tfrac{1}{2}u=\text{Const.}
\]
